% ---------------------------------------------------------------------------
\section{Conclusion}
\label{sec:conclusion}

We presented \sys{}, an agentic VLM for panoramic dental radiograph diagnosis built on a
principle we argue is the right contract between a reasoning policy and its tools in any domain
with learned components: \emph{grounding is a hypothesis to interrogate, not an answer to
trust}. Concretely, the policy learns to correct its grounding tool (\texttt{nudge\_crop},
box-relative and data-only), and it can only learn this because the training signal deliberately
exposes tool error: a tiered perturbation mechanism that is decoupled from the detector's real
accuracy --- so demonstrations never go stale --- and that never announces errors verbally ---
so the learned behavior survives to inference time. The same philosophy appears in the reward:
complexity-scaled efficiency that exempts exactly the locate-and-correct calls the task
genuinely needs, and a ground-truth-anchored, multi-finding accuracy term that an LLM judge
cannot be hacked into.

The completed evidence supports each design decision: zero-shot frontier VLMs fail on DENTEX
for structural reasons that map one-to-one onto the tool suite; the grounding tool is strong
enough to be useful and imperfect enough to need correcting; and the perturbation mechanism
elicits a clean, monotone correction gradient in 3{,}694 verified demonstrations (38.4\% of
large offsets corrected, 1.1\% of clean boxes) with recoveries up to IoU 0.066 $\to$ 0.938 ---
demonstrated judgment, not scripted steps. The end-to-end training results on the frozen
evaluation protocol are the immediate next step; the pipeline, corpus, audit tooling, and
protocol are released to make every number in this paper regenerable (artifact map and
regeneration commands in \cref{app:repro}).
